% Part: turing-machines
% Chapter: machines-computations
% Section: unary-numbers

\documentclass[../../../include/open-logic-section]{subfiles}

\begin{document}

\olfileid[fr]{tur}{mac}{una}
\olsection{Représentation unaire des nombres}

\begin{explain}
Les machines de Turing travaillent sur des suites de symboles
inscrites sur leur ruban. Selon l'alphabet qu'emploie une machine,
ces suites peuvent représenter diverses entrées et sorties. Les
machines qui calculent des fonctions \emph{arithmétiques},
c'est-à-dire des fonctions de nombres naturels, nous intéressent
particulièrement. Une façon simple de représenter ces nombres
consiste à les coder par des suites d'un seul symbole,~$\TMstroke$.
Pour $n\in\Nat$, notons $\TMstroke^n$ la suite vide si $n=0$ et,
sinon, la suite composée d'exactement $n$ symboles $\TMstroke$.
\end{explain}

\begin{defn}[Calcul]
Une machine de Turing~$M$ \emph{calcule} la fonction
$f\colon \Nat^k \to \Nat$ si, pour tous $n_1,\dots,n_k\in\Nat$,
$M$ s'arrête sur l'entrée
\[
\TMstroke^{n_1} \TMblank \TMstroke^{n_2} \TMblank \dots \TMblank \TMstroke^{n_k}
\]
avec pour sortie $\TMstroke^{f(n_1, \dots, n_k)}$.
\end{defn}

\begin{prob}
Définissez ce que signifie, pour une machine de Turing~$M$,
calculer une fonction $f\colon \Nat^k \to \Nat^m$.
\end{prob}

\begin{ex}\ollabel{ex:adder}
\emph{Addition} : construisons une machine qui calcule la fonction
$f(n,m)=n+m$. Elle doit commencer avec deux blocs de symboles
$\TMstroke$ de longueurs $n$ et~$m$ sur le ruban, puis s'arrêter
avec un seul bloc de $n+m$ symboles $\TMstroke$. Les deux blocs
de l'entrée sont séparés par un~$\TMblank$ : une méthode consiste
à écrire un trait dans cette case, puis à effacer le dernier
symbole~$\TMstroke$.\footnote{Correction éditoriale : dans le diagramme source, la lecture d'un trait dans $q_0$ mène à $q_1$. Elle doit rester dans $q_0$ pour atteindre et remplacer le séparateur ; nous corrigeons cette boucle.}
\begin{figure}\[
\begin{tikzpicture}[->,>=stealth',shorten >=1pt,auto,node distance=2.8cm,
                    semithick]
  \tikzstyle{every state}=[fill=none,draw=black,text=black]

  \node[initial,state]         (A)              {$q_0$};
  \node[state]         (B) [right of=A] {$q_1$};
  \node[state]         (C) [right of=B] {$q_2$};

  \path (A) edge node {\TMtrans{\TMblank}{\TMstroke}{\TMstay}} (B)
           edge [loop above] node {\TMtrans{\TMstroke}{\TMstroke}{\TMright}} (A)
        (B) edge [loop above] node {\TMtrans{\TMstroke}{\TMstroke}{\TMright}} (B)
            edge node {\TMtrans{\TMblank}{\TMblank}{\TMleft}} (C)
        (C) edge [loop above] node {\TMtrans{\TMstroke}{\TMblank}{\TMstay}} (C);
\end{tikzpicture}
\]\caption{Machine calculant $f(x,y)=x+y$}
\ollabel{fig:adder}
\end{figure}
\end{ex}

\begin{prob}
Suivez les configurations successives de la machine de
l'\olref[tur][mac][una]{ex:adder} pour l'entrée~$\tuple{3,2}$.
Que se passe-t-il si elle calcule $0+0$ ?
\end{prob}

\begin{explain}
Dans l'\olref[rep]{ex:doubler}, nous avons présenté une machine
de Turing qui reçoit une suite de symboles $\TMstroke$ et
s'arrête avec deux fois plus de ces symboles sur le ruban :
la machine doubleuse. Toutefois, sa sortie comporte des cases
blanches à gauche du bloc de traits doublé. Contrairement à ce
qu'on pourrait croire, elle ne calcule donc pas la fonction
$f(x)=2x$ au sens de la définition précédente. Nous allons
décrire deux façons de remédier à ce problème.
\end{explain}

\begin{ex}
La machine de la \olref{fig:doubler-disc} calcule la fonction
$f(x)=2x$. Au lieu d'effacer l'entrée et d'écrire deux traits
tout à droite pour chaque trait lu dans l'entrée, comme le fait
la machine de l'\olref[rep]{ex:doubler}, elle ajoute un seul trait
à droite pour chacun de ces traits. Elle doit repérer où se
termine l'entrée : elle laisse donc une case blanche entre
l'entrée et les traits ajoutés, qu'elle remplit d'un trait à
la toute fin. Pour « mémoriser » sa position dans l'entrée,
elle remplace temporairement un trait du bloc initial par
une case blanche.
  \begin{figure}
\[
\begin{tikzpicture}[->,>=stealth',shorten >=1pt,auto,node distance=2.8cm,
                    semithick]
  \tikzstyle{every state}=[fill=none,draw=black,text=black]

  \node[initial,state] (0)              {$q_0$};
  \node[state]         (1) [above of=0] {$q_1$};
  \node[state]         (2) [above of=1] {$q_2$};
  \node[state]         (3) [right of=2] {$q_3$};
  \node[state]         (4) [below of=3] {$q_4$};
  \node[state]         (5) [below of=4] {$q_5$};
  \node[state]         (6) [above right of=4] {$q_6$};
  \node[state]         (7) [below of=6] {$q_7$};
  \node[state]         (8) [below of=7] {$q_8$};

  \path (0) edge node {\TMtrans{\TMstroke}{\TMblank}{\TMright}} (1)
%    (0) edge node {\TMtrans{\TMblank}{\TMblank}{\TMstay}} (h)
    (1) edge [loop left] node {\TMtrans{\TMstroke}{\TMstroke}{\TMright}} (1)
      edge node {\TMtrans{\TMblank}{\TMblank}{\TMright}} (2)
    (2) edge [loop above] node {\TMtrans{\TMstroke}{\TMstroke}{\TMright}} (2)
        edge node {\TMtrans{\TMblank}{\TMstroke}{\TMleft}} (3)
    (3) edge [loop above] node {\TMtrans{\TMstroke}{\TMstroke}{\TMleft}} (3)
        edge node[left] {\TMtrans{\TMblank}{\TMblank}{\TMleft}} (4)
    (4) edge        node[left] {\TMtrans{\TMstroke}{\TMstroke}{\TMleft}} (5)
    (5) edge [loop below] node {\TMtrans{\TMstroke}{\TMstroke}{\TMleft}} (5)
        edge              node {\TMtrans{\TMblank}{\TMstroke}{\TMright}} (0)
    (4) edge      node[sloped] {\TMtrans{\TMblank}{\TMstroke}{\TMright}} (6)
    (6) edge              node {\TMtrans{\TMblank}{\TMstroke}{\TMright}} (7)
    (7) edge [loop right] node {\TMtrans{\TMstroke}{\TMstroke}{\TMright}} (7)
        edge              node {\TMtrans{\TMblank}{\TMblank}{\TMleft}} (8)
    (8) edge [loop right] node {\TMtrans{\TMstroke}{\TMblank}{\TMstay}} (8);
    \end{tikzpicture}
\]
\caption{Machine calculant $f(x)=2x$}
\ollabel{fig:doubler-disc}
\end{figure}
\end{ex}

\begin{ex}\ollabel{ex:mover}
Une autre façon de calculer $f(x)=2x$ consiste à conserver la
première machine doubleuse et à lui ajouter, à la fin, des
états et des instructions qui déplacent le bloc doublé de
traits tout à gauche du ruban. La machine de la \olref{fig:mover}
accomplit cette seule dernière tâche : si elle démarre sur un
ruban comportant un bloc de cases blanches suivi d'un bloc
de traits, la tête étant placée n'importe où dans le bloc
de cases blanches, elle efface les traits un à un et les écrit
au début du ruban. Pour reconnaître qu'elle a terminé, elle
marque d'abord la fin du bloc de traits avec un symbole
$\TMendtape$, qu'elle efface ensuite. Nous avons commencé
à numéroter ses états à~$q_6$ pour pouvoir les adjoindre à
ceux de la machine doubleuse. Il suffit d'ajouter l'instruction
$\delta(q_0,\TMblank)=\tuple{q_6,\TMblank,\TMstay}$,
autrement dit une flèche de~$q_0$ vers~$q_6$ étiquetée
$\TMtrans{\TMblank}{\TMblank}{\TMstay}$. Une difficulté
subsiste : la machine ainsi obtenue ne fonctionne pas pour
l'entrée~$x=0$. Nous en faisons un exercice.
\begin{figure}
  \[
  \begin{tikzpicture}[->,>=stealth',shorten >=1pt,auto,node distance=2.8cm,
                      semithick]
    \tikzstyle{every state}=[fill=none,draw=black,text=black]
    \node[initial,state] (6)              {$q_6$};
    \node[state]         (7) [right of=6] {$q_7$};
    \node[state]         (8) [right of=7] {$q_8$};
    \node[state]         (9) [below of=8] {$q_9$};
    \node[state]         (10) [left of=9] {$q_{10}$};
    \node[state]         (11) [left of=10]  {$q_{11}$};
    \node[state]         (12) [below of=11]  {$q_{12}$};
    \node[state]         (13) [right of=12] {$q_{13}$};
    \node[state]         (14) [right of=13] {$q_{14}$};
    \path
    (6)  edge node {\TMtrans{\TMblank}{\TMblank}{\TMright}} (7)
    (7)  edge [loop above] node {\TMtrans{\TMblank}{\TMblank}{\TMright}} (7)
         edge node {\TMtrans{\TMstroke}{\TMstroke}{\TMright}} (8)
    (8)  edge [loop above] node {\TMtrans{\TMstroke}{\TMstroke}{\TMright}} (8)
         edge node {\TMtrans{\TMblank}{\TMendtape}{\TMleft}} (9)
    (9)  edge [loop right] node {\TMtrans{\TMstroke}{\TMstroke}{\TMleft}} (9)
         edge node {\TMtrans{\TMblank}{\TMblank}{\TMright}} (10)
    (10) edge node[above] {\TMtrans{\TMstroke}{\TMblank}{\TMleft}} (11)
    (11) edge [loop above] node {\TMtrans{\TMblank}{\TMblank}{\TMleft}} (11)
         edge node[left] {\begin{tabular}{@{}l@{}}
          \TMtrans{\TMendtape}{\TMendtape}{\TMright}\\
          \TMtrans{\TMstroke}{\TMstroke}{\TMright}
         \end{tabular}} (12)
    (12) edge node[] {\TMtrans{\TMblank}{\TMstroke}{\TMright}} (13)
    (13) edge [loop below] node {\TMtrans{\TMblank}{\TMblank}{\TMright}} (13)
         edge node[sloped] {\TMtrans{\TMstroke}{\TMblank}{\TMleft}} (11)
    (13) edge node {\TMtrans{\TMendtape}{\TMblank}{\TMstay}} (14);
  \end{tikzpicture}
  \]
  \caption{Déplacement d'un bloc de symboles $\TMstroke$ vers la gauche}
  \ollabel{fig:mover}
\end{figure}
\end{ex}

\begin{prob}
Dans l'\olref[tur][mac][una]{ex:mover}, nous avons décrit une
machine obtenue en combinant la machine doubleuse de la
\olref[tur][mac][rep]{fig:doubler} avec la machine qui
déplace les traits de la \olref[tur][mac][una]{fig:mover}.
Que se passe-t-il si cette machine combinée reçoit
l'entrée~$x=0$, c'est-à-dire un ruban vide ? Comment la
modifier pour qu'elle s'arrête alors avec la sortie~$2x=0$ ?
Vous pouvez le faire en ajoutant un état et
une transition, et en modifiant la transition qui appelle la machine de déplacement.\footnote{Corrections éditoriales : l'exercice source renvoie à la doubleuse déjà normalisée, plutôt qu'à la première doubleuse, et annonce seulement des ajouts. Le renvoi est rétabli ; une transition existante peut aussi être modifiée pour distinguer l'entrée vide avant le déplacement.}
\end{prob}

\begin{prob}
\emph{Soustraction} : concevez une machine de Turing qui,
sur une entrée formée de deux suites non vides de traits,
de longueurs $n$ et~$m$ avec $n>m$, calcule la fonction
$f(n,m)=n-m$.
\end{prob}

\begin{prob}
\emph{Égalité} : concevez une machine de Turing qui
calcule la fonction suivante :
\[
\fn{equality}(n,m) = 
\begin{cases}
  \text{1} & \text{si~$n = m$} \\
  \text{0} & \text{si~$n \neq m$}
\end{cases}
\]
où $n$ et~$m$ appartiennent à~$\PosInt$.
\end{prob}

\begin{prob}
Concevez une machine de Turing qui calcule
$\min(x,y)$ pour des entiers positifs $x$ et~$y$
représentés sur le ruban par deux suites de traits
séparées par un~$\TMblank$. Vous pouvez employer
d'autres symboles dans l'alphabet de la machine.

La fonction $\min$ retient le plus petit de ses
arguments : $\min(3,5)=3$, $\min(20,16)=16$ et
$\min(4,4)=4$, par exemple.
\end{prob}

\begin{defn}
Une machine de Turing~$M$ calcule la fonction
partielle $f\colon\Nat^k\pto\Nat$ si, pour toute
entrée de la forme
$\TMstroke^{n_1}\concat\TMblank\concat
\dots\concat\TMblank\concat\TMstroke^{n_k}$ :
\begin{enumerate}
\item si $f(n_1,\dots,n_k)=m$, alors $M$ s'arrête
  avec la sortie $\TMstroke^m$;
\item si $f(n_1,\dots,n_k)$ n'est pas définie,
  alors $M$ ne s'arrête pas, ou bien elle s'arrête
  avec une sortie qui n'est pas un seul bloc de
  symboles $\TMstroke$.
\end{enumerate}
La chaîne vide représente le bloc de zéro traits : une sortie vide signifie donc la valeur $0$, et ne peut servir de marqueur pour une valeur indéfinie. \end{defn}

\end{document}
